\section*{Resumen}

Colombia enfrenta una crisis crítica en el sector salud derivada de la deuda acumulada de aproximadamente 32.9 billones de pesos entre 29 Entidades Promotoras de Salud (EPS), según informes del Ministerio de Salud emitidos en julio de 2025. Esta problemática financiera ha desencadenado graves consecuencias para las Instituciones Prestadoras de Servicios de Salud (IPS), generando desabastecimiento de medicamentos, demoras extensas en la entrega de fármacos y barreras de acceso para pacientes que dependen de tratamientos continuos. La Corte Constitucional ha constatado que persisten dificultades en la entrega de medicamentos por deudas acumuladas, incremento de tutelas y ausencia de información confiable para alertar sobre desabastecimientos.

Ante esta situación, el presente trabajo propone el desarrollo de un modelo de inteligencia artificial que brinde una solución temporal mientras se implementan reformas estructurales al sistema de salud. La propuesta consiste en un sistema integrado que combina Redes Neuronales Convolucionales (CNN) con el algoritmo K-Nearest Neighbors (KNN) para clasificar medicamentos a partir de imágenes y recomendar alternativas farmacéuticas similares. El modelo utiliza aprendizaje supervisado donde la CNN actúa como extractor de características, generando embeddings que permiten al KNN identificar medicamentos con componentes activos similares o presentaciones alternativas.

La solución implementa un pipeline de machine learning 100\% automatizado que, a partir de una fotografía del medicamento, clasifica su tipo (antibiótico, antiinflamatorio o analgésico), extrae información relevante del fármaco y recomienda alternativas disponibles basándose en similitud de componentes. Este enfoque busca empoderar a los pacientes con información oportuna mientras se resuelven los problemas estructurales del sistema de salud colombiano, mitigando el impacto de la crisis sobre personas con enfermedades crónicas o tratamientos de alto costo.

\subsection*{Palabras clave}
Medicamentos, EPS, IPS, Redes Neuronales, Algoritmo, Regresión, Clasificación, Modelo, Machine Learning

\section*{Abstract}

Colombia faces a critical crisis in the health sector derived from the accumulated debt of approximately 32.9 trillion pesos among 29 Health Promoting Entities (EPS), according to reports from the Ministry of Health issued in July 2025. This financial problem has triggered severe consequences for Health Service Provider Institutions (IPS), generating medicine shortages, extensive delays in drug delivery, and access barriers for patients who depend on continuous treatments. The Constitutional Court has confirmed that difficulties persist in medicine delivery due to accumulated debts, increased legal actions, and absence of reliable information to alert about shortages.

Facing this situation, the present work proposes the development of an artificial intelligence model that provides a temporary solution while structural reforms to the health system are implemented. The proposal consists of an integrated system that combines Convolutional Neural Networks (CNN) with the K-Nearest Neighbors (KNN) algorithm to classify medicines from images and recommend similar pharmaceutical alternatives. The model uses supervised learning where the CNN acts as a feature extractor, generating embeddings that allow the KNN to identify medicines with similar active components or alternative presentations.

The solution implements a 100\% automated machine learning pipeline that, from a photograph of the medicine, classifies its type (antibiotic, anti-inflammatory, or analgesic), extracts relevant information about the drug, and recommends available alternatives based on component similarity. This approach seeks to empower patients with timely information while structural problems of the Colombian health system are resolved, mitigating the crisis impact on people with chronic diseases or high-cost treatments.

\subsection*{Key words}
Medicines, EPS, IPS, Neural Networks, Algorithm, Regression, Classification, Modeling, Machine Learning

\pagebreak